\chapter{Formalisation}

\section{Parsing goal and modelling stance}
The intended task is the following: given a finite corpus $\mathcal C$ of sentences and a string of
words drawn from the corpus vocabulary, we want to assign Lambek types to that string and,
when several parses are available, rank them probabilistically. In other words, the model is not
only a recogniser of grammaticality but a \emph{ranker} of candidate analyses.

The formalisation in this thesis deliberately simplifies several practical issues.
\begin{itemize}
    \item It assumes that lexical type probabilities are already available, either from manual
    annotation or from a prior learning stage.
    \item It does not handle words outside the corpus vocabulary.
    \item It treats probabilities as corpus-relative confidence values rather than as a fully learned
    generative model.
    \item It abstracts away from implementation issues such as chart parsing, estimation, and
    smoothing.
\end{itemize}
These choices are not defects so much as scope restrictions: they isolate the semantic core that
is needed for the soundness proof.


\section{Syntax of the associative Lambek calculus}
Let $\At$ be a non-empty set of atomic types. The set of Lambek formulae is generated by the
grammar
\[
A ::= p \mid A \backslash A \mid A / A \mid A \concat A
\qquad (p \in \At).
\]
We work with the associative Lambek calculus, written $\LC$, in an algebraic presentation. The
primitive relation $\vdash_{\LC}$ is a preorder on formulae generated by the following validities:
\begin{align*}
& A \vdash_{\LC} A, \\
& (A \concat B) \concat C \vdash_{\LC} A \concat (B \concat C), \\
& A \concat (B \concat C) \vdash_{\LC} (A \concat B) \concat C, \\
& A \concat B \vdash_{\LC} C \iff A \vdash_{\LC} C/B, \\
& A \concat B \vdash_{\LC} C \iff B \vdash_{\LC} A\backslash C, \\
& A \vdash_{\LC} B \text{ and } B \vdash_{\LC} C \Rightarrow A \vdash_{\LC} C,
\end{align*}

This is equivalent in spirit to the usual residuated presentation of the associative Lambek
calculus. It is particularly convenient here because the soundness proof proceeds by checking the
validity of identity, associativity, residuation, and transitivity under the proposed semantics.


\section{Corpus fragments and vocabulary}
Let $\mathcal C$ be a finite corpus of sentences over a finite vocabulary $\Sigma$. We denote $\Sigma^+$ to be the set of finite non-empty strings from vocabulary $\Sigma$. In order to keep the
semantics finite and directly tied to parsing data, let $\Frag(\mathcal C)$ be the set of all non-empty
contiguous strings that occur in at least one sentence in $\mathcal C$. Thus:
\[
\Frag(\mathcal C) \subseteq \Sigma^+.
\]
Because $\mathcal C$ is finite, the set $\Frag(\mathcal C)$ is finite as well.

The choice of $\Frag(\mathcal C)$ rather than all of $\Sigma^+$ is deliberate. It reflects the fact that
this semantics is trained and evaluated relative to observed data. It also ensures that the minima
and maxima appearing in the recursive clauses below are taken over finite sets. Extending the
semantics to arbitrary strings over $\Sigma$ would require a separate treatment of unseen data,
which is left to future work.


\section{Lexical probabilities}
We assume a lexical probability assignment
\[
\Lex : \Sigma \times \At \to [0,1]
\]
such that for every word $w \in \Sigma$,
\[
\sum_{p \in \At} \Lex(w,p) = 1,
\]
where all but finitely many summands are zero. Intuitively, $\Lex(w,p)$ is the probability that
the word $w$ is assigned the atomic type $p$.

\begin{remark}
If a word $w \in \Sigma$ occurs in the corpus but is never assigned an atomic type $p$, then
$\Lex(w,p)=0$. This matches the intuition from the draft: lack of attestation for a type within
the observed vocabulary should count against that type.
\end{remark}

\begin{remark}
Words outside the corpus vocabulary are not assigned probabilities here at all. Rather than
assigning them an arbitrary value such as $0$ or leaving them as exceptional elements with
undefined behaviour, they are simply excluded from the domain of the semantics. This keeps the
formal system clean and makes explicit that out-of-vocabulary handling is a future extension.
\end{remark}


\section{Probabilistic valuation}
We now define a valuation
\[
v : \Frag(\mathcal C) \times \mathrm{Fm}(\LC) \to [0,1].
\]
The definition proceeds by recursion on formula complexity.

\begin{definition}[Atomic clauses]
For $x \in \Frag(\mathcal C)$ and $p \in \At$, define
\[
v(x,p) =
\begin{cases}
\Lex(x,p) & \text{if $x$ is a single word in $\Sigma$,}\\
0 & \text{if $x$ has length greater than $1$.}
\end{cases}
\]
\end{definition}

The second clause is a modelling choice: atomic categories are taken to be lexical categories of
single words. One could allow atomic categories for longer fragments as well, but that would
require extra data and is unnecessary for the present soundness result.

\begin{definition}[Product]
For $x \in \Frag(\mathcal C)$,
\[
v(x,A\concat B)
= \max\{v(y,A)\,v(z,B) \mid y,z \in \Frag(\mathcal C),\ yz=x\},
\]
with the convention that the maximum of the empty set is $0$.
\end{definition}

This says that the probability of $x$ having composite type $A\concat B$ is the best score among
all observed binary decompositions of $x$ into a left part of type $A$ and a right part of type
$B$.

\begin{definition}[Left residual]
For $x \in \Frag(\mathcal C)$,
\[
v(x,A\backslash B)
= \min\Bigl(\{1\} \cup
\Bigl\{\min\bigl(1,\frac{v(w,B)}{v(y,A)}\bigr)
\ \Bigm|\ y \in \Frag(\mathcal C),\ yx=w \in \Frag(\mathcal C),\ v(y,A)>0\Bigr\}\Bigr).
\]
\end{definition}

\begin{definition}[Right residual]
For $x \in \Frag(\mathcal C)$,
\[
v(x,B/A)
= \min\Bigl(\{1\} \cup
\Bigl\{\min\bigl(1,\frac{v(w,B)}{v(y,A)}\bigr)
\ \Bigm|\ y \in \Frag(\mathcal C),\ xy=w \in \Frag(\mathcal C),\ v(y,A)>0\Bigr\}\Bigr).
\]
\end{definition}


\section{Why these clauses are well-defined}
The original draft raised two technical problems: division by zero in the residual definitions and
the treatment of words not present in the corpus. The present formalisation resolves both.

First, division by zero is avoided by restricting the residual minima to witnesses with positive
denominator, i.e. to contexts $y$ such that $v(y,A)>0$. If no such witness exists, the set in the
minimum collapses to $\{1\}$ and the value becomes $1$. This is a natural default: in the absence
of any positively supported counterevidence, the residual connective is not forced below $1$.
The inner truncation $\min(1,\_)$ ensures that the value stays in the interval $[0,1]$.

Second, out-of-vocabulary words are excluded from the domain entirely. The semantics is not a
universal language model over all possible strings; it is a finite, corpus-relative semantics for
observed fragments. This is consistent with the stated goal of the thesis, namely to parse strings
of words drawn from the corpus vocabulary.

\begin{remark}[Subtleties not captured by the present formalisation]
The following issues are intentionally left outside the formal model.
\begin{enumerate}[label=(\alph*)]
    \item \textbf{Learning from raw corpora.} The formalisation assumes lexical probabilities are
    already given, whereas a practical system would need to infer them.
    \item \textbf{Out-of-vocabulary generalisation.} Smoothing, backoff, or subword modelling would be
    needed for real applications.
    \item \textbf{Global normalisation.} The values $v(x,A)$ behave as probabilities locally, but the model is
    not presented here as a full generative probability distribution over derivation trees.
    \item \textbf{Unseen but grammatical strings.} By working over $\Frag(\mathcal C)$, the semantics is tied to
    attested fragments rather than all strings over $\Sigma$.
\end{enumerate}
These restrictions should be read as part of the thesis scope, not as permanent limitations on
future versions of the model.
\end{remark}


\section{Semantic consequence}
\begin{definition}
For formulae $A$ and $B$, define
\[
A \Sem B
\quad\text{iff}\quad
v(x,A) \leq v(x,B) \text{ for all } x \in \Frag(\mathcal C).
\]
\end{definition}

Thus $A \Sem B$ means that every observed fragment is at least as probable as a realisation of
$B$ as it is as a realisation of $A$. This is the semantic preorder with respect to which soundness
will be proved.

